\section{Extensions}
\label{sec:extensions}

Recent studies have made significant strides in extending the capabilities of MLLMs, spanning from more potent foundational abilities to broader coverage of scenarios. We trace the principal development of MLLMs in this regard.

\noindent \textbf{Granularity Support.}
To facilitate better interaction between agents and users, researchers have developed MLLMs with finer support of granularities in terms of model inputs and outputs.
On the input side, models that support finer control from user prompts are developed progressively, evolving from image to region~\cite{chen2023shikra,zhang2023gpt4roi,xuan2023pink} and even pixels~\cite{yuan2023osprey,rasheed2023glamm,you2023ferret}.
Specifically, Shikra~\cite{chen2023shikra} supports region-level input and understanding. Users may interact with the assistant more flexibly by referring to specific regions, which are represented in bounding boxes of natural language forms.
Ferret~\cite{you2023ferret} takes a step further and supports more flexible referring by devising a hybrid representation scheme. The model supports different forms of prompts, including point, box, and sketch. 
Similarly, Osprey~\cite{yuan2023osprey} supports point input by utilizing a segmentation model~\cite{sam}. Aided by the exceptional capabilities of the pre-trained segmentation model, Osprey enables specifying a single entity or part of it with a single click.
On the output side, grounding capabilities are improved in line with the development of input support. 
Shikra~\cite{chen2023shikra} supports response grounded in the image with box annotations, resulting in higher precision and finer referring experience.
LISA~\cite{lai2023lisa} further supports mask-level understanding and reasoning, which makes pixel-level grounding possible. 

\noindent \textbf{Modality Support.}
Increased support for modalities is a tendency for MLLM studies.
On the one hand, researchers have explored adapting MLLMs to support the input of more multimodal content, such as 3D point cloud~\cite{xu2023pointllm,chen2023ll3da,huang2023embodied,hong20233d}.
On the other hand, MLLMs are also extended to generate responses of more modalities, such as image~\cite{emu,zhan2024anygpt,wu2023next,aiello2023jointly}, audio~\cite{zhang2023speechgpt,rubenstein2023audiopalm,zhan2024anygpt,wu2023next}, and video~\cite{wu2023next,wang2024modaverse}. 
For example, NExT-GPT~\cite{wu2023next} proposes a framework that supports inputs and outputs of mixed modalities, specifically, combinations of text, image, audio, and video, with the help of diffusion models~\cite{ho2020denoising,rombach2022high} attached to the MLLM. The framework applies an encoder-decoder architecture and puts LLM as a pivot for understanding and reasoning.

\noindent \textbf{Language Support.}
Current models are predominantly unilingual, probably due to the fact that high-quality non-English training corpus is scarce. Some works have been devoted to developing multilingual models so that a broader range of users can be covered.
VisCPM~\cite{hu2023large} transfers model capabilities to the multilingual setting by designing a multi-stage training scheme. Specifically, the scheme takes English as a pivotal language, with abundant training corpus. Utilizing a pre-trained bilingual LLM, the multimodal capabilities are transferred to Chinese by adding some translated samples during instruction tuning. 
Taking a similar approach, Qwen-VL~\cite{bai2023qwen} is developed from the bilingual LLM Qwen~\cite{qwen-lm} and supports both Chinese and English. During pre-training, Chinese data is mixed into the training corpus to preserve the bilingual capabilities of the model, taking up 22.7\% of the whole data volume.


\noindent \textbf{Scenario/Task Extension.}
Apart from developing common general-purpose assistants, some studies have focused on more specific scenarios where practical conditions should be considered, while others extend MLLMs to downstream tasks with specific expertise.

A typical tendency is to adapt MLLMs to more specific real-life scenarios. MobileVLM~\cite{chu2023mobilevlm} explores developing small-size variants of MLLMs for resource-limited scenarios. Some designs and techniques are utilized for deployment on mobile devices, such as LLMs of smaller size and quantization techniques to speed up computation.
Other works develop agents that interact with real-world~\cite{gong2023mindagent,huang2023embodied,embodiedgpt}, \eg user-friendly assistants specially designed for Graphical User Interface (GUI), as exemplified by CogAgent~\cite{hong2023cogagent}, AppAgent~\cite{yang2023appagent}, and Mobile-Agent~\cite{wang2024mobile}. These assistants excel in planning and guiding through each step to fulfill a task specified by users, acting as helpful agents for human-machine interaction.
Another line is to augment MLLMs with specific skills for solving tasks in different domains, \eg document understanding~\cite{ye2023mplug,liu2024textmonkey,hu2024mplug,ye2023ureader} and medical domains~\cite{llava-med,pmc-vqa,moor2023med}.
For document understanding, mPLUG-DocOwl~\cite{ye2023mplug} utilizes various forms of document-level data for tuning, resulting in an enhanced model in OCR-free document understanding. TextMonkey~\cite{liu2024textmonkey} incorporates multiple tasks related to document understanding to improve model performance. Apart from conventional document image and scene text datasets, position-related tasks are added to reduce hallucinations and help models learn to ground responses in the visual information.
MLLMs can also be extended to medical domains by instilling knowledge of the medical domain. For example, LLaVA-Med~\cite{li2023llava} injects medical knowledge into vanilla LLaVA~\cite{llava} and develops an assistant specialized in medical image understanding and question answering.
